\section{Discussion}
\label{sec:discussion}

\paragraph{Did \tool find a new bug?} No, and we now have a precise reason
why, not just an impression. Two sites on current master have the same
shape as CVE-2022-4450 (\texttt{X509\_PUBKEY\_set},
\texttt{ossl\_ml\_kem\_i2d\_pubkey}; Observation~\ref{obs:rq1-mastertriage}).
Re-checking both against the real pinned source (\texttt{results/\allowbreak
reachability\_rq1\_candidates.md}) found that our original reachability note
for \texttt{X509\_PUBKEY\_set} was itself wrong: it blamed a lock/atomic-op
failure in \texttt{EVP\_PKEY\_up\_ref}, but every platform variant of
\texttt{CRYPTO\_UP\_REF} in this codebase always returns \texttt{true} --
that failure mode is dead code. The real trigger for both sites is the same
in shape: the tracked object's state (a reference count, a key's public-key
field) must already have changed between two reads that a single
straight-line call provably treats as equal. Only a second thread
racing the same object, unsynchronized, can cause that change. We checked every
in-tree caller of both functions and found none that creates this
precondition internally; it requires a pre-existing bug or data race one
level up the call stack. We did not build a multithreaded fault-injection
test to force that race window, so we still report both as candidates,
not findings, but on a corrected and sharper basis than before.

\paragraph{Is this just matching the one patch it was built from?} The
rule (Section~\ref{sec:rule}) names no OpenSSL function. A second,
unrelated fix, \texttt{d2217c88df}, is detected without changing the
rule (Observation~\ref{obs:rq1-conformance}), and a sibling site that fix
missed is flagged and later fixed independently
(Observation~\ref{obs:rq3-timetravel}). This is evidence the rule is more
than a rewrite of one patch, but it is two data points, not a rate
(Section~\ref{sec:threats}). Separately, is the release-without-reset
\emph{convention} itself a recurring problem in this codebase, independent
of our tool? A commit-message search of \texttt{openssl/openssl}'s real
history (\texttt{results/historical\_recurrence\_search.md}) turned up ten
further, independently-authored fixes of this exact shape. They span
nine files in \texttt{crypto/}, \texttt{providers/}, and \texttt{ssl/},
from 2023 to six weeks before our pinned snapshot. None is a parameter or
BIO-chain shape either engine currently tracks (mostly struct-field and
bare local-variable resets), so this corroborates the convention's
generality without being a tool finding.

\paragraph{Why build a custom engine instead of using an existing
framework such as CodeQL or Semgrep directly?} Table~\ref{tab:rq2-comparison}
is exactly that attempt, run with each tool's stock configuration. None
of three general-purpose open-source tools distinguishes the vulnerable
revision from the fixed one for any of the three CVEs we checked
(Observation~\ref{obs:rq2-generic-miss}); their rule and query sets simply
do not encode this project's specific conventions. We instead write the
rule once, as data, over our own control-flow graph
(Section~\ref{sec:rule}), so we can inspect exactly how it was evaluated.

\paragraph{What about false positives?} In the manual triage behind
Observation~\ref{obs:rq1-mastertriage}, 1 of 18 flagged sites across all
six snapshots is a confirmed false positive: \texttt{load\_common}, which
releases only a reference-counted handle's structural reference, not the
underlying object (Section~\ref{sec:threats}). Two further sites are
suppressed by the tool itself, correctly, under the deref-after-release
rule.

\paragraph{Is a clean run on CVE-2023-0215 or CVE-2024-4741 evidence that
\tool handles them?} Mixed. Protocol B implements neither mechanism
(Section~\ref{sec:subjects}); the vulnerable/fixed distinction for those
two CVEs in Table~\ref{tab:rq2-comparison} comes from the Part~1 pattern
baseline, not \tool. Protocol A (\rqref{5}, Section~\ref{sec:rq5})
implements CVE-2023-0215's BIO-chain invariant and distinguishes its own
motivating pair. But its whole-tree run finds 0 confirmed recurrences of
11 flags, all false positives from missing ownership tracking -- so a
clean run elsewhere would not be strong evidence either, for the same
reason. CVE-2024-4741 needs a state-guard rule neither shape expresses;
future work.

\paragraph{Is the whole-tree conformance rate (91.7\% on master) a
precision number?} No. It is the fraction of release sites that are
conforming under one specific, name-based release heuristic
(Section~\ref{sec:rq1}); a site this heuristic cannot see, such as a
release behind a struct member, is absent from both the numerator and the
denominator. We leave measuring that blind spot for future work.
